\documentclass[11pt]{article}
\usepackage[a4paper,margin=1in]{geometry}
\usepackage{amsmath,amssymb,mathtools,bm}
\usepackage{enumitem,booktabs,array}
\usepackage{tcolorbox}
\usepackage{hyperref}
\usepackage[protrusion=true,expansion=false]{microtype}
\hypersetup{colorlinks=true,linkcolor=blue,urlcolor=blue,citecolor=blue}
\setlist[itemize]{topsep=2pt,itemsep=2pt}
\setlist[enumerate]{topsep=2pt,itemsep=2pt}
\newtcolorbox{keybox}{colback=blue!3!white,colframe=blue!40!black,boxrule=0.4pt,arc=1mm}
\newtcolorbox{alertbox}{colback=red!3!white,colframe=red!40!black,boxrule=0.4pt,arc=1mm}
\newtcolorbox{answerbox}{colback=green!3!white,colframe=green!40!black,boxrule=0.4pt,arc=1mm}
\newtcolorbox{drillbox}{colback=yellow!5!white,colframe=orange!60!black,boxrule=0.4pt,arc=1mm}
\newcommand{\EV}{\mathbb{E}}
\newcommand{\Var}{\mathrm{Var}}
\newcommand{\Cov}{\mathrm{Cov}}
\newcommand{\blank}[1]{\underline{\hspace{#1}}}

\title{E10-03: Howell (2017, AER) \\
\large ``Financing Innovation: Evidence from R\&D Grants'' --- Active Recall Workbook}
\author{Empirical Corporate Finance II --- Exam Prep}
\date{\today}
\begin{document}\maketitle

\begin{keybox}
\textbf{Paper in one sentence.} Using a regression-discontinuity around SBIR-grant award thresholds, Howell shows that US DOE Small Business Innovation Research (SBIR) Phase I grants \emph{cause} a $\sim$30\% increase in patenting and $\sim$20--60\% increase in follow-on private investment for small R\&D firms, with effects concentrated among young, financially-constrained firms --- government grants \emph{do} alleviate financial frictions in innovation.

\textbf{Placement.} Landmark quasi-experimental paper on public R\&D grants and innovation; anchor for innovation-finance research and innovation-policy debates. Strong identification via RD design.
\end{keybox}

\section*{Round 1: Complete Walk-Through}

\subsection*{1.1 Research question}
Do public R\&D grants generate additional innovation that would not otherwise occur, or do they crowd out private investment? Most previous work lacked clean identification.

\subsection*{1.2 Setting: SBIR grants}
\begin{itemize}
\item US DOE Small Business Innovation Research program.
\item Phase I: up to \$150K feasibility grants; Phase II: up to \$1M development grants.
\item Applications ranked by DOE review panel on technical merit; threshold cutoff funds top applicants.
\item Sharp RDD: firms just above cutoff get grant, firms just below do not.
\end{itemize}

\subsection*{1.3 Data}
\begin{itemize}
\item DOE SBIR applications 1983--2013.
\item Patents from USPTO; follow-on VC from VentureSource.
\item Firm survival, revenue growth, acquisitions.
\end{itemize}

\subsection*{1.4 Empirical design}
Regression discontinuity in the scored-rank at the award cutoff:
\[
Y_i=\alpha+\beta\cdot\mathbb{1}(\text{Rank}_i>c)+f(\text{Rank}_i-c)+\varepsilon_i,
\]
with standard RD bandwidth choice and local linear estimates.

\subsection*{1.5 Headline results}
\begin{itemize}
\item Phase I award raises patent counts by $\sim$30\%.
\item Phase I award raises follow-on VC funding probability by 10--30 pp.
\item Phase II effects smaller; diminishing returns at larger grant sizes.
\item Effects concentrated in young, financially-constrained firms --- no effect for well-funded or mature firms.
\item Consistent with relaxing financial constraints, not pure substitution.
\end{itemize}

\subsection*{1.6 Mechanism}
\begin{itemize}
\item Phase I reduces a binding constraint at the prototype / feasibility stage.
\item Grant creates credible signal that reduces private-capital screening frictions.
\item Phase II crowds out private investment partly because selected firms have already found backers.
\end{itemize}

\subsection*{1.7 Implications}
\begin{itemize}
\item Government R\&D grants can generate real innovation at the margin.
\item Program design matters: small early-stage grants more effective than larger later-stage.
\item Provides an identification benchmark for other innovation-policy research.
\end{itemize}

\subsection*{1.8 Caveats}
\begin{alertbox}
\begin{itemize}
\item Generalizability beyond DOE-SBIR questionable.
\item External validity limited to small, innovative firms.
\item RD design assumes continuity at cutoff --- strategic application-gaming could bias.
\end{itemize}
\end{alertbox}

\section*{Round 2: Level 1 Fill-in}
\begin{enumerate}
\item Setting: DOE \blank{2cm} program.
\item Identification: \blank{2cm} discontinuity.
\item Phase I patent effect: $+$\blank{1cm}\%.
\item VC follow-on effect: $+$\blank{1cm}--\blank{1cm} pp.
\item Effect concentrated in \blank{2cm}-constrained firms.
\end{enumerate}

\section*{Round 3: Level 2 Prompts}
\begin{enumerate}
\item Describe the SBIR program and RD design.
\item Report main findings on patents and follow-on capital.
\item Discuss mechanism: financial-constraint relaxation vs.\ signaling.
\item Compare Phase I and Phase II.
\item Policy implications.
\end{enumerate}

\section*{Round 4: Minimal Scaffolding}
From a blank page: (i) setting, (ii) RD design, (iii) results, (iv) mechanism, (v) policy.

\section*{Round 5: Blank Page}
Essay: can public R\&D grants cause innovation? Evidence from SBIR.

\vspace{6pt}
\begin{drillbox}
\textbf{Speed Drill.} (1) Program. (2) Identification. (3) Patent effect. (4) VC effect. (5) Who benefits.
\end{drillbox}

\section*{Quick Reference \& Answer Key}
\begin{answerbox}
\begin{itemize}
\item \textbf{Program.} DOE SBIR Phase I/II.
\item \textbf{Design.} Rank-based RD.
\item \textbf{Phase I.} Patents $+30\%$, VC $+10$--$30$ pp.
\item \textbf{Lesson.} Grants relax constraints for young firms.
\end{itemize}
\end{answerbox}

\begin{keybox}
\textbf{30-second exam pitch.} Howell (2017) uses a regression-discontinuity design around SBIR Phase I/II grant award cutoffs in DOE applications 1983--2013 to identify causal effects of public R\&D funding. Phase I grants (up to \$150K) raise recipient firms' patent counts by $\sim$30\% and follow-on VC probability by 10--30 pp; Phase II effects are smaller, indicating diminishing returns. Benefits concentrate in young, financially-constrained firms consistent with relaxing financial constraints at the feasibility stage, rather than pure substitution. The paper delivers a clean quasi-experimental estimate of grant-caused innovation and informs program design: small, early-stage grants are more effective than large, late-stage subsidies. It anchors the empirical innovation-finance literature and has become the canonical RD-based innovation-policy evaluation.
\end{keybox}

\end{document}
